\documentclass[twoside, 10pt]{article}
\usepackage{geometry}
\geometry{bottom=4em,top=6em, inner=2.2cm, outer=4em,  headheight=\paperheight}
\usepackage[export]{adjustbox}
\usepackage{array}
\usepackage{amsmath}
\usepackage{amsfonts}
\usepackage{fancyhdr}
\usepackage{luacode}
\pagestyle{fancy}
\fancyhf{}
\lhead{Precalculus - BASE}
\chead{Linear v.s. Logarithmic }
\directlua{pageNum = 1}
\rhead{Practice, Page \directlua{tex.print(pageNum); pageNum = pageNum + 1; if pageNum > 3 then pageNum = 1 end}}
\usepackage{lastpage}
\usepackage{xcolor}
\usepackage{enumitem}
\usepackage{pifont}
\usepackage{graphicx}
\graphicspath{{../img}}
\usepackage{pgfplots}
\pgfplotsset{compat=1.18}
\usepackage{tabularx}
\usepackage{polynom}

\newcommand{\R}{\mathbb R}
\newcommand{\e}{{\rm e}}
\newcommand{\pobr}[1]{\left\langle#1\right\rangle}
\newcommand{\norm}[1]{\lVert #1 \rVert}
\newcommand{\abs}[1]{\lvert #1 \rvert}

\DeclareMathOperator{\xd}{d\!}
\DeclareMathOperator{\proj}{proj}

\title{}
\date{}

\newcommand{\template}{
\directlua{tex.print(utils.genVersion())}
\noindent
{\large
First Name \rule{6em}{.1pt}\hspace{\stretch{1}} Last Name \rule{6em}{.1pt}\hspace{\stretch{1}} Date \rule{1.5em}{.1pt} -- \rule{1.5em}{.1pt} -- \rule{1.5em}{.1pt}\hspace{\stretch{1}} Period \rule{2em}{.1pt}\hspace{\stretch{1}} Score \rule{2em}{.1pt}
}
\vspace{1em}
}

\begin{luacode*}
utils = require("../utils")
\end{luacode*}

\begin{document}
\template

Consider the following sequence of 10 numbers:
\[
3,\; 7,\; 12,\; 18,\; 25,\; 31,\; 42,\; 56,\; 63,\; 78
\]

Answer the questions below.

\begin{enumerate}[label=\arabic*.]
    \item Design an algorithm to find the number 58. Clearly describe each step of your algorithm, and justify that the algorithm always terminates, regardless of whether 58 is in the sequence.
\vspace{\stretch{1}}
    \item How many iterations does your algorithm take before it ends?  
    Explain how you determined this number.
\vspace{\stretch{1}}
    \item Suppose the length of the sequence is \(n\) instead of 10.  
    As \(n\) increases, how does the number of iterations needed by your algorithm change?  
    Describe the relationship using words and, if possible, mathematical notation.
    \vspace{\stretch{1}}
\end{enumerate}
\clearpage
Here is the same sequence of numbers:
\[
3,\; 7,\; 12,\; 18,\; 25,\; 31,\; 42,\; 56,\; 63,\; 78
\]

Notice that the sequence is ordered from smallest to largest. Because of this, we can apply the \emph{binary search} algorithm.

\subsection*{Description of Binary Search}

Binary search works as follows:
\begin{enumerate}[label = (\roman*)]
    \item Look at the middle number of the sequence.
    \item If the middle number is the target number, the search stops.
    \item If the target number is smaller than the middle number, repeat the process using only the left half of the sequence.
    \item If the target number is larger than the middle number, repeat the process using only the right half of the sequence.
    \item Continue this process until the number is found or there are no numbers left to check.
\end{enumerate}

Answer the questions below.

\begin{enumerate}[label=\arabic*.]
    \item Use the binary search algorithm to search for the number \(58\).  
    Show the result of each iteration, including which numbers remain under consideration at each step.
    \vspace{\stretch{1}}
    \item How many iterations does it take to determine whether \(58\) is in the sequence?
    \vspace{\stretch{1}}
    \item Suppose the length of the sequence is \(n\).  
    As \(n\) increases, how does the number of iterations required by binary search change?  
    Describe the relationship using words and, if possible, mathematical notation.
        \vspace{\stretch{1}}
\end{enumerate}
\end{document}